% Proposition K: collapsed configurations, the critical leak, and the orientation of the orbit as a design input
% (PARADIGM.md, section 6e).
\section{Proposition K: collapsed configurations, the critical leak, and the orientation of the orbit as a design input}\label{sec:thmK}

Branch 2 of Section~\ref{sec:program} asks whether the circulation of the reference orbit can be designed. Theorem J
(Section~\ref{sec:thmJ}) makes the question precise for the collapsed configuration, the trap that Theorem H's dither
was introduced to escape: its stability is decided by the eigenvalues of a Hessian that the orientation shifts at
order $k_e^{-2}$.

\paragraph{Setting.} As in Theorem J. A configuration is \emph{collapsed} if $\vartheta_i = \bar\vartheta$ for all
$i$. Then $\fhat_\Theta = \tilde\Phi(\cdot\,;\bar\vartheta)$, so $\Theta$ is a critical point of $J_m$ if and only if
$\bar\vartheta$ is a critical point of the single-kernel objective
$J_1(\vartheta) := \frac12\int\norm{f - \tilde\Phi(\cdot\,;\vartheta)}^2\dd\rho
+ \frac{\sigma_0}{2}\norm{\vartheta - \vartheta^0}^2_{\Ginv}$. Write $\bar r := f - \tilde\Phi(\cdot\,;\bar\vartheta)$
and
\begin{gather*}
  M_2 := \int\nabla_\vartheta\tilde\Phi(\cdot\,;\bar\vartheta)^\T\nabla_\vartheta\tilde\Phi(\cdot\,;\bar\vartheta)\dd\rho\succeq0,\\
  R_2 := \int\sum_k\bar r_k\nabla^2_\vartheta\tilde\Phi_k(\cdot\,;\bar\vartheta)\dd\rho
  = \nabla^2_\vartheta\bigl\langle\bar r^\T\tilde\Phi(\cdot\,;\vartheta)\bigr\rangle_\rho\Big|_{\bar\vartheta}
  \quad\text{(residual frozen)}.
\end{gather*}
$R_2$ is the Hessian of $-V_L[\bar\mu](\vartheta)$ at the atom: the second-order cost of moving a small amount of mass
off the atom at frozen residual.

\begin{thm}{Proposition}{K}[collapsed configurations, the critical leak, and the orientation as a design input]
Let $\Theta^*$ be a collapsed critical point of $J_m$ with atom $\bar\vartheta$.
\begin{enumerate}[label=\textup{(\roman*)}]
  \item \emph{Symmetry splitting.}
    $\nabla^2J_m(\Theta^*) = I_m\otimes(\sigma_0\Ginv - R_2) + \frac1m(\mathbf 1\mathbf 1^\T)\otimes M_2$. On the
    separation subspace $\{(u_1,\dots,u_m) : \sum_iu_i = 0\}$, of dimension $(m-1)p$, it acts as
    $\sigma_0\Ginv - R_2$; on the collective subspace $\{u_i\text{ all equal}\}$ as
    $\nabla^2J_1(\bar\vartheta) = \sigma_0\Ginv - R_2 + M_2$. The same splitting holds for the Jacobian of any
    permutation-equivariant particle drift at a collapsed configuration, in particular for the second-order averaged
    drift of Remark H.7a \textup{(}Section~\ref{sec:finitegain}\textup{)}: it has the form
    $I_m\otimes D + \frac1m(\mathbf 1\mathbf 1^\T)\otimes E$.
  \item \emph{Separation stability at leading order.} $\Theta^*$ is a non-degenerate local minimum of $J_m$ iff
    $\nabla^2J_1(\bar\vartheta)\succ0$ and $\sigma_0\Ginv\succ R_2$, that is, iff $\bar\vartheta$ is a non-degenerate
    minimiser of $J_1$ and $\sigma_0 > \sigma_{\mathrm{crit}}(\bar\vartheta) := \lambda_{\max}(G^{1/2}R_2G^{1/2})$.
    The separation Floquet exponents of the collapsed periodic orbit of \textup{(CL)} are then, by Theorem J(b),
    $(1/k_e)\lambda_j(\sigma_0I - G^{1/2}R_2G^{1/2}) + O(k_e^{-2})$, each with multiplicity $m-1$, and they do not
    depend on $m$. The condition $\sigma_0\Ginv\succ R_2$ is the statement that the atom is a strict local minimiser
    of the first-variation potential $V_L[\bar\mu] + \frac{\sigma_0}{2}\norm{\cdot}^2_{\Ginv}$ on its support, the
    local form of the Chizat--Bach optimality condition for measures.
  \item \emph{The orientation-odd shift.} Let $\mathcal O$ be traversed in either direction \textup{(}the same $\rho$,
    hence the same $\Theta^*$ and the same leading exponents\textup{)}. Write the second-order normal form of
    Theorem J's proof as
    \begin{equation}
      \dot\Theta = -\frac1{k_e}G_m\nabla J_m(\Theta) + b(\Theta) + O(k_e^{-3}),\qquad
      b := b_{\mathrm{frozen}} + b_{\mathrm{osc}} = O(k_e^{-2}),
      \tag{N}
    \end{equation}
    with $b_{\mathrm{frozen}} = \frac1{k_e^2}G_m\nabla(U_{\mathrm{even}} + U_{\mathrm{odd}})$ \textup{(}Proposition
    H.7\textup{)} and $b_{\mathrm{osc}}$ of Remark H.7a. \textup{(N)} has a unique equilibrium $\bar\Theta_2$ near
    $\Theta^*$, which is collapsed, with $\bar\Theta_2 - \Theta^* = -J_1^{-1}b(\Theta^*) + O(k_e^{-2})$ for
    $J_1 := -\frac1{k_e}G_m\nabla^2J_m(\Theta^*)$ \textup{(}the first-order location of Theorem J(d)\textup{)}; let
    $J_2$ be the Jacobian of \textup{(N)} at $\bar\Theta_2$. The parameter Floquet exponents of the collapsed orbit
    are the eigenvalues of $-J_2$ up to $O(k_e^{-3})$, and
    \[
      J_2 = J_1 + \partial b(\Theta^*) + \partial J_1(\Theta^*)[\bar\Theta_2 - \Theta^*] + O(k_e^{-3}),
    \]
    where the third term --- the third derivative of $J_m$ contracted with the first-order offset --- is $O(k_e^{-2})$,
    of the same order as $\partial b$, and must be kept: the Jacobian of \textup{(N)} at $\Theta^*$ is not the right
    linearisation. Under orbit reversal $\rho$, $\Theta^*$, $J_1$ and $U_{\mathrm{even}}$ are unchanged while
    $U_{\mathrm{odd}}$ and $b_{\mathrm{osc}}$ change sign \textup{(}Proposition H.7(d), Remark H.7a\textup{)}, so
    with $b_{\mathrm{even}} := \frac1{k_e^2}G_m\nabla U_{\mathrm{even}}$ and
    $b_{\mathrm{odd}} := \frac1{k_e^2}G_m\nabla U_{\mathrm{odd}} + b_{\mathrm{osc}}$,
    \begin{gather*}
      J_2 = J_1 + E \pm O,\\
      E := \partial b_{\mathrm{even}}(\Theta^*) - \partial J_1(\Theta^*)\bigl[J_1^{-1}b_{\mathrm{even}}(\Theta^*)\bigr],\qquad
      O := \partial b_{\mathrm{odd}}(\Theta^*) - \partial J_1(\Theta^*)\bigl[J_1^{-1}b_{\mathrm{odd}}(\Theta^*)\bigr],
    \end{gather*}
    for the two orientations. $J_1$, $E$ and $O$ are permutation-equivariant and preserve the separation subspace;
    by \textup{(i)} their restrictions to one separation block are $p\times p$ matrices, written
    $J_1|_{\mathrm{sep}} = -\frac1{k_e}(\sigma_0I - GR_2)$, $D_{\mathrm{even}}$, $D_{\mathrm{odd}}$. The separation
    Floquet exponents of the two orientations differ by
    \[
      \lambda_j^{(+)} - \lambda_j^{(-)} = -2\,\frac{w_j^\T D_{\mathrm{odd}}v_j}{w_j^\T v_j} + O(k_e^{-3})
    \]
    \textup{(}simple eigenvalue $\lambda_j$ of $-J_1|_{\mathrm{sep}}$, right and left eigenvectors $v_j$,
    $w_j$\textup{)}, an $O(k_e^{-2})$ quantity that scales like $k_e^{-2}$ at fixed $(G,\sigma_0)$, while the
    exponents themselves are $O(1/k_e)$. The difference is a property of the excitation's orientation alone: the
    occupation measure, the averaged system $\Aavg$ and all of Theorems B--H are blind to it.
  \item \emph{The critical leak as a design point.} At $\sigma_0 = \sigma_{\mathrm{crit}}(\bar\vartheta)$ the softest
    separation exponent vanishes at leading order and equals the $k_e^{-2}$ value
    $-(w^\T D_{\mathrm{even}}v)/(w^\T v) \mp (w^\T D_{\mathrm{odd}}v)/(w^\T v) + O(k_e^{-3})$. If
    $\abs{w^\T D_{\mathrm{odd}}v} > \abs{w^\T D_{\mathrm{even}}v}$, the collapsed configuration is stable for one
    orientation and unstable for the other: the direction in which the reference is traversed decides whether the
    kernels separate. More generally, for $\sigma_0$ in a window of width $O(1/k_e)$ around $\sigma_{\mathrm{crit}}$
    --- the window where the leading exponent $(\sigma_0 - \sigma_{\mathrm{crit}})(1 + o(1))/k_e$ lies between the two
    $k_e^{-2}$ values --- the orientation decides the sign, and the centre of the window,
    $\sigma_0 = \sigma_{\mathrm{crit}} + k_e(w^\T D_{\mathrm{even}}v)/(w^\T v) + O(k_e^{-1})$, can be computed from
    the averaged data alone.
\end{enumerate}
\end{thm}
\status{(i) and (ii) are proved (linear algebra at the collapsed point and Theorem J(b)). (iii) and (iv) are proved
given the second-order normal form (N) with an $O(k_e^{-3})$ remainder, which is Theorem J's proof carried one order
further (periodic averaging to second order, Sanders--Verhulst--Murdock~\cite[Ch.~2--3]{svm2007}; the first-order step
is written out in the proof of Theorem J, the second-order step is Remark H.7a's derivation and is not written out
here). The orientation-odd shift, its $k_e^{-2}$ scaling, the prediction and the design window are checked numerically
in one instance (\code{exp_orientation_stability.py}).}

\begin{proof}
\textit{(i)} $\partial^2J_m/\partial\vartheta_i\partial\vartheta_j
= \frac1m\int\nabla\tilde\Phi_i^\T\nabla\tilde\Phi_j\dd\rho - \delta_{ij}\int\sum_kr_k\nabla^2\tilde\Phi_{i,k}\dd\rho
+ \delta_{ij}\sigma_0\Ginv$ (the first term from the quadratic dependence of $\frac12\int\norm{r}^2\dd\rho$ on the
$m$ features with weights $1/m$, the second from the $\vartheta_i$-curvature of the residual term:
with $r = f - \frac1m\sum\tilde\Phi_j$,
$\partial_{\vartheta_i}(m\bar L_m) = -\int r^\T\nabla\tilde\Phi_i\dd\rho$ and
$\partial_{\vartheta_j}\partial_{\vartheta_i}(m\bar L_m)
= \frac1m\int\nabla\tilde\Phi_j^\T\nabla\tilde\Phi_i\dd\rho - \delta_{ij}\int\sum_kr_k\nabla^2\tilde\Phi_{i,k}\dd\rho$).
At a collapsed point all blocks coincide, giving the Kronecker form; the two subspaces are the isotypic components of
the permutation action, on which $\frac1m(\mathbf 1\mathbf 1^\T)$ acts as $0$ and $1$ respectively. The second-order
averaged drift is permutation-equivariant (each particle moves by a function of its own position and of the measure),
so its Jacobian at a collapsed point has the same form with $D$ the self-derivative and $E$ the derivative through
the measure.

\smallskip
\textit{(ii)} Positive definiteness of a block-diagonal form; the exponents are Theorem J(b) applied to the separation
block, whose eigenvalues of $G_m^{1/2}\nabla^2J_mG_m^{1/2}$ restricted to the separation subspace are those of
$\sigma_0I - G^{1/2}R_2G^{1/2}$, each repeated $m-1$ times. For the last sentence, the first variation of
$\mathcal G := \bar L + \frac{\sigma_0}{2}\int\norm{\cdot - \vartheta^0}^2_{\Ginv}\dd\mu$ at
$\bar\mu = \delta_{\bar\vartheta}$ is
$V_L[\bar\mu](\vartheta) + \frac{\sigma_0}{2}\norm{\vartheta - \vartheta^0}^2_{\Ginv}
= -\bigl\langle\bar r^\T\tilde\Phi(\cdot\,;\vartheta)\bigr\rangle_\rho + \frac{\sigma_0}{2}\norm{\vartheta - \vartheta^0}^2_{\Ginv}$,
whose Hessian at $\bar\vartheta$ is $-R_2 + \sigma_0\Ginv$.

\smallskip
\textit{(iii)} The proof of Theorem J carries one order further: the near-identity change of variables of its Step~3
removes the oscillating terms through order $k_e^{-2}$, which turns \textup{(CL)} into \textup{(N)} with an
$O(k_e^{-3})$ remainder (this is the normal form of Remark H.7a's derivation, with $b_{\mathrm{osc}}$ the term that
the frozen-measure expansion misses). \textup{(N)} is a regular perturbation of the averaged system near the
hyperbolic point $\Theta^*$, so it has a unique equilibrium $\bar\Theta_2$ within $O(1/k_e)$ of $\Theta^*$, obtained
by the implicit function theorem; by the first-order expansion,
$\bar\Theta_2 - \Theta^* = -J_1^{-1}b(\Theta^*) + O(k_e^{-2})$ ($J_1^{-1} = O(k_e)$, $b = O(k_e^{-2})$). \textup{(N)}
is permutation-equivariant and $\Theta^*$ is collapsed, so the equivariant Newton iterates stay collapsed and
$\bar\Theta_2$ is collapsed. Theorem J(b) applied to \textup{(N)} (same proof, remainder $O(k_e^{-3})$) gives the
Floquet exponents as the eigenvalues of $-J_2 + O(k_e^{-3})$. Expanding the Jacobian of \textup{(N)} about $\Theta^*$:
$J_2 = J_1 + \partial b(\Theta^*) + \partial J_1(\Theta^*)[\bar\Theta_2 - \Theta^*]
+ O(\norm{\bar\Theta_2 - \Theta^*}^2/k_e) + O(k_e^{-2}\norm{\bar\Theta_2 - \Theta^*})$, and both error terms are
$O(k_e^{-3})$; substituting the offset gives $E$ and $O$, which are the even and odd parts under reversal because
$\partial J_1$ is even and $b_{\mathrm{even}}$, $b_{\mathrm{odd}}$ have the parities named. All three operators are
permutation-equivariant, hence preserve the separation subspace, and by the Kronecker structure of \textup{(i)} each
acts on it as $I_{m-1}\otimes(\text{a $p\times p$ block})$. First-order eigenvalue perturbation of the simple
eigenvalue $\lambda_j$ of $-J_1|_{\mathrm{sep}}$ gives
\[
  \lambda_j^{(\pm)} = \lambda_j - \frac{w_j^\T D_{\mathrm{even}}v_j}{w_j^\T v_j}
  \mp \frac{w_j^\T D_{\mathrm{odd}}v_j}{w_j^\T v_j} + O(k_e^{-3}),
\]
and the difference as stated. (For the $(m-1)$-fold degenerate separation modes the degeneracy is the trivial one of
the Kronecker structure, which the restriction to one $p$-dimensional block removes; a further degeneracy within the
block would call for degenerate perturbation theory.)

\smallskip
\textit{(iv)} is \textup{(iii)} at a vanishing leading eigenvalue; the window follows because the leading exponent is
$(\sigma_0 - \sigma_{\mathrm{crit}})/k_e\cdot(1 + o(1))$ near $\sigma_{\mathrm{crit}}$ ($\bar\vartheta$, $R_2$ and
the eigenvectors depend smoothly on $\sigma_0$) while the shifts are $O(k_e^{-2})$; its centre is where the leading
exponent cancels the even shift.
\end{proof}

\begin{rmk}{Remark}{K.1}[reading]
Three things.
\begin{itemize}
  \item The collapsed configuration's stability is a single-kernel question at leading order (the Hessian
    $\sigma_0\Ginv - R_2$ does not depend on $m$), which explains why the $m = 6$ and $m = 32$ experiments of
    Section~\ref{sec:thmH} collapsed alike: a leak above $\sigma_{\mathrm{crit}}$ traps every network size.
  \item Dither (Theorem H) and the orientation (this proposition) are two different ways to leave the trap: the first
    adds entropy, the second tilts the Hessian by $O(k_e^{-2})$. The second is free but small; it matters only near
    $\sigma_{\mathrm{crit}}$, that is, when the leak is tuned to the edge of trapping, or for nearly flat separation
    modes such as the weak saddle of Remark H.2.
  \item The sign of the shift depends on the data $(f,\mathcal O,\bar\vartheta)$ through $D_{\mathrm{odd}}$ and has
    no general sign: the design rule is to compute $D_{\mathrm{odd}}$ (from the averaged data: $U_{\mathrm{odd}}$,
    $b_{\mathrm{osc}}$ and the third derivative of $J_m$) and choose the orientation with the smaller softest
    separation exponent (or the negative one); the window of (iv) is where the choice matters.
\end{itemize}
\end{rmk}

\paragraph{Numerics} (\code{exp_orientation_stability.py} $\to$ \code{results/orientation_stability/table.md}).
Instance A of Theorem J's numerics (Section~\ref{sec:thmJ}; $n = 2$, $m = 3$, the elliptic reference traversed in both
directions, $G = \operatorname{diag}(4,4,1,1,1)$). For each orientation the collapsed periodic orbit is found by
Newton shooting (residual $\le4\cdot10^{-13}$), the parameter block of its monodromy matrix is restricted to the
separation and collective subspaces, and the exponents are compared with the leading order
$(\sigma_0I - G^{1/2}R_2G^{1/2})/k_e$ and with the second-order prediction, the eigenvalues of $-J_2$ at the
equilibrium $\bar\Theta_2$ of (N) (Newton from $\Theta^*$; every $\bar\Theta_2$ found is collapsed to $10^{-8}$). The
second-order equilibrium is $\norm{\bar\Theta_2 - \Theta^*}_{\Ginv} = 3.7\cdot10^{-2}$, $1.8\cdot10^{-2}$,
$9.0\cdot10^{-3}$, $4.5\cdot10^{-3}$ away from $\Theta^*$ at $k_e = 10$, $20$, $40$, $80$: the $O(1/k_e)$ offset of
Theorem J(d).
\begin{itemize}
  \item \emph{The shift and its prediction} ($\sigma_0 = 0.5$, a stable collapse). The softest separation exponents
    of the two orientations differ by $9.0\cdot10^{-4}$, $1.9\cdot10^{-4}$, $4.1\cdot10^{-5}$, $9.7\cdot10^{-6}$ at
    $k_e = 10$, $20$, $40$, $80$: fitted exponent $-2.18$ against $k_e$ ((iii): $-2$). Scaled by $k_e^2$, the
    measured corrections to the leading order are $-0.032$, $-0.057$, $-0.070$, $-0.077$ for $(+)$ and $-0.123$,
    $-0.132$, $-0.136$, $-0.139$ for $(-)$, approaching the predicted $-0.083$ and $-0.143$ with a gap that halves at
    every doubling of $k_e$ (the $O(k_e^{-3})$ remainder); the predicted difference at $\bar\Theta_2$ ---
    $7.8\cdot10^{-4}$, $1.7\cdot10^{-4}$, $3.9\cdot10^{-5}$, $9.4\cdot10^{-6}$ --- is within 14\%, 9\%, 5\%, 3\% of
    the measured one. The prediction with the Jacobian taken at $\Theta^*$ instead of $\bar\Theta_2$ (the
    third-derivative term omitted) reads $9.0\cdot10^{-4}$, $2.1\cdot10^{-4}$, $5.0\cdot10^{-5}$, $1.2\cdot10^{-5}$:
    it agrees at $k_e = 10$ by accident and its $k_e^2$-scaled value converges to $0.078$ while the measured one
    converges to $0.060$, the 30\% error that the numerics exposed (\code{review/REVIEW_LOG.md}, round 5, K-1). The
    separation exponents come in exact pairs (multiplicity $m - 1 = 2$) and the collective exponents differ between
    the orientations by 2\% at $k_e = 10$ and 0.2\% at $k_e = 80$, as (i)--(iii) describe.
  \item \emph{At the critical leak} ($\sigma_0 = \sigma_{\mathrm{crit}} = 0.4606$). The leading exponent vanishes and
    the $k_e^{-2}$ terms decide: the softest separation exponent is $-4.1\cdot10^{-4}$ for $(+)$ and
    $-1.27\cdot10^{-3}$ for $(-)$ at $k_e = 10$, $-4.8\cdot10^{-5}$ and $-8.7\cdot10^{-5}$ at $k_e = 40$ --- unstable
    for both orientations in this instance, the even part dominating the odd one by a factor 2--3; the predictions at
    $\bar\Theta_2$ are within 14\% and 6\% of the measured values at $k_e = 40$.
  \item \emph{The orientation decides} ($\sigma_0$ designed from the averaged data). Solving
    $\mathrm{pred}_+(\sigma_0) + \mathrm{pred}_-(\sigma_0) = 0$ for the predicted softest exponents gives
    $\sigma_0 - \sigma_{\mathrm{crit}} = 0.0125$ at $k_e = 10$ and $0.0061$ at $k_e = 20$, the centre of the window of
    (iv); there the measured softest separation exponents are $+8.5\cdot10^{-4}$ and $-2.5\cdot10^{-5}$ ($k_e = 10$),
    $+1.4\cdot10^{-4}$ and $-3.7\cdot10^{-5}$ ($k_e = 20$): the same reference orbit, traversed one way, keeps the
    collapsed kernels together, and traversed the other way makes them separate. The window is narrow --- the leading
    exponent at these $\sigma_0$ is $1.2\cdot10^{-3}$ and $3.0\cdot10^{-4}$ --- and the design lands inside it
    although the $k_e^{-3}$ terms move both measured exponents by about $+4.7\cdot10^{-4}$ and $+6\cdot10^{-5}$
    against the prediction.
\end{itemize}
\checkedby{\code{tests/test_periodic.py}: the Hessian splitting (i) against finite differences ($10^{-5}$ relative) at
the collapsed minimiser of instance A, the isotypic eigenvalue structure, the gap equation of $\sigma_{\mathrm{crit}}$
($10^{-8}$), and the two-fold zero of the separation block at $\sigma_{\mathrm{crit}}$ with a positive definite
collective block (\code{test_proposition_K_hessian_splitting_and_critical_leak}); the experiment above for
(iii)--(iv) (\code{exp_orientation_stability.py} $\to$ \code{results/orientation_stability/table.md}).}
